\documentclass[12pt]{article}
\usepackage[margin=0.75in]{geometry}

\newcommand{\duedate}{04/05/2026}
\newcommand{\assignment}{When Do Curricula Work: Summary}

\newcommand{\name}{Aksel Kretsinger-Walters, adk2164}
\newcommand{\email}{adk2164@columbia.edu}

\makeatletter
\def\input@path{{../}{../../}{../../../}}
\makeatother
\input{pset_template_CLMM.tex} %% DO NOT CHANGE THIS LINE

\begin{document}
\psetheader %% DO NOT CHANGE THIS LINE

\section*{When Do Curricula Work?}

\paragraph{Motivation.}
Curriculum learning --- presenting training examples in a meaningful order,
typically easy-to-hard --- has been proposed as a way to improve convergence
and generalization. Anti-curriculum learning (hard-to-easy) has also shown
benefits in some settings. Despite a rich literature, empirical results
conflict: curricula help in some experiments but not others, and no ordered
learning method consistently improves across contexts. This paper aims to
systematically determine when and why curricula are beneficial.

\paragraph{Implicit Curricula.}
Before studying explicit orderings, the authors investigate whether neural
networks already learn examples in a consistent order under standard i.i.d.\
training. They define the \emph{learned iteration} of a sample as the epoch
at which the model first correctly and permanently classifies it.

Across 142 configurations (architectures $\times$ optimizers) on CIFAR-10,
the learned iteration ordering is highly consistent within architecture
families. The Spearman correlation between orderings from different CNNs
is very high, indicating that a robust, architecture-dependent notion of
example difficulty exists. This \emph{implicit curriculum} emerges from the
interaction of architecture and optimization, not from explicit data ordering.

\paragraph{Curriculum Framework.}
The paper formalizes curricula via three components:

\begin{enumerate}
		\item \textbf{Scoring function} $s(\mathbf{x}) \in \mathbb{R}$: assigns
				a difficulty score to each example. Three are studied: loss of a
				reference model, learned iteration, and c-score (consistency of
				correct predictions across independent training runs).
		\item \textbf{Pacing function} $g(t)$: determines how many examples are
				available at training step $t$, controlling how quickly the full
				dataset is introduced. Six families are considered: logarithmic,
				exponential, step, linear, quadratic, and root.
		\item \textbf{Order}: curriculum (easy-to-hard), anti-curriculum
				(hard-to-easy), or random ordering within the paced subset.
\end{enumerate}

All three scoring functions are highly correlated, so c-score is used
throughout.

\paragraph{Main Finding: Curricula Provide Marginal Benefit in Standard Settings.}
Across a massive sweep of over 25{,}000 models on CIFAR-10, CIFAR-100,
FOOD101, and FOOD101N, the paper finds:

\begin{itemize}
		\item Under standard training (full budget, clean data), \textbf{no ordering
						--- curriculum, anti-curriculum, or random --- provides statistically
				significant improvement} over standard i.i.d.\ SGD.
		\item Any marginal gains from pacing functions are explained entirely by the
				\emph{dynamic dataset size} effect (starting with a smaller dataset
				and growing it), not by the ordering itself. Random-curriculum performs
				as well as curriculum.
\end{itemize}

\paragraph{When Curricula Do Help.}
Curricula become genuinely beneficial in two specific regimes:

\begin{itemize}
		\item \textbf{Limited training time:} When the compute budget is
				drastically reduced (e.g., 352 steps instead of 17{,}600),
				curriculum ordering provides a clear advantage over anti-curriculum
				and standard training. At the tightest budgets, even random-curriculum
				(random order with a pacing function) outperforms the baseline,
				suggesting the dynamic dataset size itself helps.
		\item \textbf{Noisy labels:} With 20--80\% label noise on CIFAR-100,
				curriculum learning significantly outperforms all other methods.
				The easy-to-hard ordering allows the model to first learn clean
				patterns before encountering corrupted examples. The best pacing
				functions under noise correspond to simply ignoring all noisy data
				during training.
\end{itemize}

These findings generalize to larger-scale datasets (FOOD101, FOOD101N).

\paragraph{Why These Regimes.}
Both limited time and label noise shift the c-score distribution toward
harder examples, making the distinction between easy and hard examples more
consequential. In standard settings with sufficient compute and clean data,
SGD eventually learns all examples regardless of order, so curricula add
nothing. When resources are constrained or data is corrupted, prioritizing
reliable examples first becomes valuable.

\paragraph{Limitations.}
The study is limited to image classification with fixed datasets (CIFAR,
FOOD101) and does not cover language modeling, where curricula are standard
practice in large-scale pretraining (GPT-3, T5). The scoring functions are
computed before training begins, so self-paced or online difficulty
estimation is not considered. The paper also notes that conclusions about
curricula may not transfer to very different task types or training regimes.

\paragraph{Takeaway.}
Curricula are not a general-purpose improvement to training. Their benefit
is concentrated in settings where learning resources are scarce (limited
compute or corrupted data), making example prioritization valuable. In
standard well-resourced training, the order of examples is largely
irrelevant --- i.i.d.\ sampling works just as well.

\end{document}
